% ============================================================
%  CONCLUSION
% ============================================================
\section{Conclusion}

\subsection{Summary of Findings}

We introduced NanoChat 2.0, a systematic study of four recent architectural and training
innovations applied to a depth-12, approximately 85-million-parameter autoregressive language
model.
Our controlled experiments on the FineWeb-Edu corpus yield the following principal conclusions:

\begin{itemize}

  \item \textbf{ReLU\(^2\) is the best activation for this scale.}
  All six activation functions tested--including evolutionary-search activations (GeLUSine, GSP)
  and custom designs (GMTU, Turbulent)--underperform the NanoChat $\mathrm{ReLU}^2$ default on
  aggregate ChatCORE.
  Benefits reported in computer vision do not transfer to small-scale language modelling.

  \item \textbf{MTP improves knowledge tasks at a throughput cost.}
  Multi-token prediction with $n=4$ heads yields modest gains on ARC and MMLU relative to the
  baseline, consistent with the hypothesis that multi-horizon objectives build richer factual
  representations.
  However, MTP incurs the worst training throughput of all variants, limiting its practical
  utility under fixed compute budgets at this scale.

  \item \textbf{Attention Residuals do not benefit shallow models.}
  The PreNorm dilution problem that AttnRes \parencite{chen2026attnres} is designed to address is
  negligible at depth 12.
  The additional optimisation complexity of learned depth-wise attention weights is not amortised
  within 5{,}000 training steps, resulting in worse validation BPB and ChatCORE than the
  standard additive baseline.

  \item \textbf{Engram is not recommended at NanoChat scale.}
  Despite matching the baseline on pretraining BPB, the Engram layer badly hurts compositional
  generation (SpellingBee, HumanEval) after SFT, while providing only marginal improvements on
  knowledge retrieval tasks.
  The mechanism's benefits were demonstrated at 27B parameters / 262B tokens--scales
  approximately two orders of magnitude larger than ours.

  \item \textbf{FP8 mixed-precision is the most impactful single change.}
  Enabling FP8 matrix multiplications consistently improves every post-SFT metric without any
  architectural modification.
  We hypothesise that FP8 quantisation noise acts as implicit regularisation in this
  data-limited regime, and recommend enabling \texttt{--fp8} by default for H100 training runs.

\end{itemize}

\subsection{Limitations}

Several limitations bound the scope of our conclusions.
First, all experiments are restricted to 5{,}000 training iterations---approximately one-tenth of
the compute used in the NanoChat leaderboard submission.
Techniques such as MTP and AttnRes, which require more data and steps to fully manifest their
representational advantages, may show different results under a longer training budget.
Second, hardware heterogeneity between experiment groups (4 $\times$ H100 vs.\
2 $\times$ H100) introduces a confounding factor in throughput comparisons, which we partially
address by reporting in GPU-hours.
Third, the hyperparameters of Engram (layer positions, memory size, gate initialisation) were not
extensively tuned for our scale, and a more careful configuration search might mitigate some of
the observed degradation.
Fourth, our analysis of FP8's beneficial effect is speculative; a proper ablation comparing equal
throughput in FP16 vs.\ FP8 would be needed to isolate the regularisation effect.

\subsection{Future Work}

Several directions emerge naturally from our findings.
\begin{itemize}
  \item \textbf{Longer training}
  Re-running MTP and AttnRes with 50{,}000+ iterations would determine whether their
  representational advantages eventually outweigh the added complexity.
  \item \textbf{Engram mitigations}
  If Engram is revisited, potential improvements include: (i) reducing the module's influence by
  down-scaling the engram weight, (ii) placing engram at a later layer (e.g.\ layer 3 instead of
  1) to interfere less with early-layer reasoning, and (iii) augmenting the SFT data mixture with
  more SpellingBee-style and code examples to teach the gate to suppress the memory contribution
  for compositional tasks.
  \item \textbf{RLHF fine-tuning}
  The best-performing configuration (FP8 baseline) could be further improved through
  reinforcement learning from human feedback (RLHF) to enhance conversational quality and
  instruction following.
  \item \textbf{Combined modifications}
  Testing the joint effect of MTP + FP8, or investigating whether Engram and MTP are
  complementary, remains an open question.
\end{itemize}

\subsection{Broader Impact}

All in all, our work demonstrates a reproducible proof-of-concept for conducting controlled experiments on GPT-2-scale chatbots on a modest academic GPU node.
By identifying which recent innovations transfer to small-scale language modelling and which do not, we provide practical guidance for researchers and practitioners seeking to deploy capable autoregressive models in resource-constrained settings--including offline use, on-device mobile
inference, and applications where the cost of a frontier-scale model is unjustifiable. This democratises empirical LLM research, empowering small teams with limited budgets to obtain reliable, reproducible results.
